\chapter{Position and Momentum Fermions}
\label{ch:fermions}

\section{Introduction and Physical Framework}

Having established the abstract Canonical Anticommutation Relations (CAR) and Fock space representations in Chapter~\ref{ch:car_fock}, and the discrete Fourier transform in Chapter~\ref{ch:fourier}, we now unite these structures to describe physical fermions on the one-dimensional lattice.

Throughout this chapter and the remainder of the monograph, the lattice length is assumed to be an even positive integer:
\begin{equation}
L = 2h, \quad h \in \mathbb{N}, \quad h \ge 1.
\end{equation}
The single-particle index set $\iota$ is instantiated as the centered momentum band:
\begin{equation}
\iota = \Lambda^* = \left\{ k \in \mathbb{Z} \;\middle|\; -L < 2k \le L \right\} = \{-h+1, -h+2, \dots, h-1, h\},
\end{equation}
endowed with the standard linear order inherited from $\mathbb{Z}$. 

The fermionic state space is the Fock space over the momentum band:
\begin{equation}
\mathcal{F}(L) \coloneqq \ell^2(\mathcal{P}(\Lambda^*)) \equiv \bigoplus_{S \subseteq \Lambda^*} \mathbb{C} |S\rangle,
\end{equation}
which has exact complex dimension $\dim_{\mathbb{C}} \mathcal{F}(L) = 2^L$. The momentum creation and annihilation operators $\{c_k, c_k^\dagger \mid k \in \Lambda^*\}$ satisfy the standard CAR algebra:
\begin{equation}
\{c_k, c_p^\dagger\} = \delta_{k, p} I, \qquad \{c_k, c_p\} = 0, \qquad \{c_k^\dagger, c_p^\dagger\} = 0, \qquad c_k^\dagger = (c_k)^\dagger.
\end{equation}

\section{Spatial Fermions via the Unitary Fourier Transform}

In continuous quantum field theory, field operators $\psi(x)$ are operator-valued distributions that must be smeared with test functions. On our discrete spatial circle $\Lambda = \mathbb{Z}/L\mathbb{Z}$, spatial fermions are well-defined bounded operators constructed as finite linear combinations of momentum modes.

Recall from Chapter~\ref{ch:fourier} the real normalization factor $a \coloneqq 1/\sqrt{L} \in \mathbb{R}_{>0}$ satisfying $L a^2 = 1$, and the canonical primitive $L$-th root $\zeta = e^{2\pi i / L}$.

\begin{definition}[Spatial Fermionic Operators]
For every spatial site $x \in \Lambda$, the spatial annihilation operator $c_x$ and spatial creation operator $c_x^\dagger$ are defined by:
\begin{align}
c_x &\coloneqq \frac{1}{\sqrt{L}} \sum_{k \in \Lambda^*} \chi(k, x)\, c_k = \frac{1}{\sqrt{L}} \sum_{k \in \Lambda^*} \zeta^{k x_{\mathrm{val}}}\, c_k, \label{eq:spatial_c} \\
c_x^\dagger &\coloneqq \frac{1}{\sqrt{L}} \sum_{k \in \Lambda^*} \chi(-k, x)\, c_k^\dagger = \frac{1}{\sqrt{L}} \sum_{k \in \Lambda^*} \zeta^{-k x_{\mathrm{val}}}\, c_k^\dagger. \label{eq:spatial_cdag}
\end{align}
The local spatial fermion number operator is defined by:
\begin{equation}
n_x \coloneqq c_x^\dagger c_x.
\end{equation}
\end{definition}

\begin{lemma}[Hilbert Adjoint Compatibility]
The spatial creation operator is the exact Hilbert adjoint of the spatial annihilation operator:
\begin{equation}
(c_x)^\dagger = c_x^\dagger \quad \text{for all } x \in \Lambda.
\end{equation}
\end{lemma}

\begin{proof}
Taking the Hilbert adjoint of the linear combination~\eqref{eq:spatial_c}:
\[
(c_x)^\dagger = \left( \frac{1}{\sqrt{L}} \sum_{k \in \Lambda^*} \zeta^{k x_{\mathrm{val}}} c_k \right)^\dagger = \frac{1}{\sqrt{L}} \sum_{k \in \Lambda^*} \overline{\zeta^{k x_{\mathrm{val}}}} (c_k)^\dagger.
\]
Because $(c_k)^\dagger = c_k^\dagger$ and $\overline{\zeta^{k x_{\mathrm{val}}}} = \zeta^{-k x_{\mathrm{val}}} = \chi(-k, x)$, this sum coincides with the definition of $c_x^\dagger$ in~\eqref{eq:spatial_cdag}.
\end{proof}

\section{Exact Position CAR Inheritance}

The fundamental theorem of discrete spatial fermions is that they inherit the local canonical anticommutation relations without distortion.

\begin{theorem}[Exact Spatial CAR]
\label{thm:spatial_car}
For all spatial sites $x, y \in \Lambda$, the spatial operators satisfy:
\begin{align}
\{c_x, c_y^\dagger\} &= \delta_{x, y} I, \label{eq:pos_car_mixed} \\
\{c_x, c_y\} &= 0, \label{eq:pos_car_ann} \\
\{c_x^\dagger, c_y^\dagger\} &= 0. \label{eq:pos_car_cre}
\end{align}
\end{theorem}

\begin{proof}
We evaluate each anticommutator using the bilinearity of $\{\cdot, \cdot\}$:
\begin{enumerate}
    \item For the mixed anticommutator:
    \begin{align*}
    \{c_x, c_y^\dagger\} &= \left\{ \frac{1}{\sqrt{L}} \sum_{k \in \Lambda^*} \chi(k, x) c_k, \; \frac{1}{\sqrt{L}} \sum_{p \in \Lambda^*} \chi(-p, y) c_p^\dagger \right\} \\
    &= \frac{1}{L} \sum_{k \in \Lambda^*} \sum_{p \in \Lambda^*} \chi(k, x) \chi(-p, y) \{c_k, c_p^\dagger\}.
    \end{align*}
    Substituting the momentum CAR $\{c_k, c_p^\dagger\} = \delta_{k, p} I$:
    \[
    \{c_x, c_y^\dagger\} = \frac{1}{L} \sum_{k \in \Lambda^*} \chi(k, x) \chi(-k, y) I = \frac{1}{L} \left( \sum_{k \in \Lambda^*} \zeta^{k (x_{\mathrm{val}} - y_{\mathrm{val}})} \right) I.
    \]
    By Theorem~\ref{thm:fourier_orthogonality}(2) (dual character orthogonality), the momentum sum equals $L\, \delta_{x, y}$. Therefore:
    \[
    \{c_x, c_y^\dagger\} = \frac{1}{L} (L\, \delta_{x, y}) I = \delta_{x, y} I.
    \]

    \item For the annihilation anticommutator:
    \[
    \{c_x, c_y\} = \frac{1}{L} \sum_{k, p \in \Lambda^*} \chi(k, x) \chi(p, y) \{c_k, c_p\} = 0,
    \]
    because $\{c_k, c_p\} = 0$ for all $k, p \in \Lambda^*$. Taking adjoints yields $\{c_x^\dagger, c_y^\dagger\} = 0$.
\end{enumerate}
\end{proof}

\begin{remark}[Formalization Note]
In the verified Lean~4 core (Chapter 5, Position and Momentum Fermions), this result is proved purely through bilinear expansion and character orthogonality. Notice an essential conceptual feature of this construction: spatial fermions are obtained directly as global Fourier transforms of momentum fermions, rather than through local Jordan--Wigner transformations from spin operators. Consequently, there are no non-local string operators connecting spatial modes to momentum modes.
\end{remark}

\section{Operator Fourier Inversion}

The relationship between position and momentum fermions is completely symmetric under Fourier inversion.

\begin{theorem}[Fourier Inversion for Operators]
\label{thm:op_fourier_inversion}
The momentum operators are recovered from spatial fermions via the inverse discrete Fourier transform:
\begin{align}
c_k &= \frac{1}{\sqrt{L}} \sum_{x \in \Lambda} \chi(-k, x)\, c_x = \frac{1}{\sqrt{L}} \sum_{x \in \Lambda} \zeta^{-k x_{\mathrm{val}}}\, c_x, \label{eq:inv_ck} \\
c_k^\dagger &= \frac{1}{\sqrt{L}} \sum_{x \in \Lambda} \chi(k, x)\, c_x^\dagger = \frac{1}{\sqrt{L}} \sum_{x \in \Lambda} \zeta^{k x_{\mathrm{val}}}\, c_x^\dagger. \label{eq:inv_ckdag}
\end{align}
\end{theorem}

\begin{proof}
Substitute the definition of $c_x$ from~\eqref{eq:spatial_c} into the right-hand side of~\eqref{eq:inv_ck}:
\[
\frac{1}{\sqrt{L}} \sum_{x \in \Lambda} \chi(-k, x) \left( \frac{1}{\sqrt{L}} \sum_{p \in \Lambda^*} \chi(p, x) c_p \right) = \frac{1}{L} \sum_{p \in \Lambda^*} \left( \sum_{x \in \Lambda} \chi(p - k, x) \right) c_p.
\]
By Theorem~\ref{thm:fourier_orthogonality}(1) (spatial character orthogonality), $\sum_{x \in \Lambda} \chi(p - k, x) = L\, \delta_{p, k}$. Thus:
\[
\frac{1}{L} \sum_{p \in \Lambda^*} (L\, \delta_{p, k}) c_p = c_k.
\]
Taking Hilbert adjoints yields the corresponding identity for $c_k^\dagger$.
\end{proof}

\section{Conservation of Total Particle Number}

\begin{theorem}[Particle Number Invariance]
\label{thm:number_invariance}
The total particle number operator $\hat{N}_{\mathrm{tot}}$ satisfies:
\begin{equation}
\hat{N}_{\mathrm{tot}} = \sum_{k \in \Lambda^*} c_k^\dagger c_k = \sum_{x \in \Lambda} c_x^\dagger c_x = \sum_{x \in \Lambda} n_x.
\end{equation}
\end{theorem}

\begin{proof}
Expand the sum of spatial number operators using~\eqref{eq:spatial_c} and~\eqref{eq:spatial_cdag}:
\begin{align*}
\sum_{x \in \Lambda} c_x^\dagger c_x &= \sum_{x \in \Lambda} \left( \frac{1}{\sqrt{L}} \sum_{k \in \Lambda^*} \chi(-k, x) c_k^\dagger \right) \left( \frac{1}{\sqrt{L}} \sum_{p \in \Lambda^*} \chi(p, x) c_p \right) \\
&= \frac{1}{L} \sum_{k, p \in \Lambda^*} c_k^\dagger c_p \left( \sum_{x \in \Lambda} \chi(p - k, x) \right).
\end{align*}
Applying spatial character orthogonality $\sum_{x \in \Lambda} \chi(p - k, x) = L\, \delta_{p, k}$:
\[
\sum_{x \in \Lambda} c_x^\dagger c_x = \frac{1}{L} \sum_{k, p \in \Lambda^*} c_k^\dagger c_p (L\, \delta_{p, k}) = \sum_{k \in \Lambda^*} c_k^\dagger c_k = \hat{N}_{\mathrm{tot}}.
\]
\end{proof}

\section{The Dirac Fermi Sea and Linearized Hamiltonians}

In relativistic field theory, the physical vacuum is not the bare state $|\emptyset\rangle$ devoid of particles, but the filled Dirac sea where all negative energy states are occupied.

\begin{definition}[Dirac Sea Configuration and Vacuum]
For even lattice length $L = 2h$, the half-filled Dirac Fermi sea configuration $S_\Omega \subset \Lambda^*$ consists of all non-positive momentum modes:
\begin{equation}
S_\Omega \coloneqq \{k \in \Lambda^* \mid k \le 0\} = \{-h+1, -h+2, \dots, -1, 0\}.
\end{equation}
The cardinality of the Fermi sea is $\#S_\Omega = h = L/2$. The Dirac sea vacuum state is the corresponding occupation state:
\begin{equation}
|\Omega\rangle \coloneqq |S_\Omega\rangle \in \mathcal{F}(L).
\end{equation}
\end{definition}

\begin{definition}[Bare and Physical Momentum Hamiltonians]
\begin{enumerate}
    \item The bare kinetic Hamiltonian $H_0$ is defined by:
    \begin{equation}
    H_0 \coloneqq \sum_{k \in \Lambda^*} k\, c_k^\dagger c_k = \sum_{k \in \Lambda^*} k\, n_k.
    \end{equation}
    \item The ground state sea energy $E_\Omega$ is the eigenvalue of $H_0$ on $|\Omega\rangle$:
    \begin{equation}
    E_\Omega \coloneqq \sum_{k \in S_\Omega} k.
    \end{equation}
    \item The physical, normal-ordered Hamiltonian $\hat{P}$ is defined by:
    \begin{equation}
    \hat{P} \coloneqq H_0 - E_\Omega \cdot I,
    \end{equation}
    so that $\hat{P} |\Omega\rangle = 0$.
\end{enumerate}
\end{definition}

\begin{lemma}[Evaluation of Ground State Energy]
For $L = 2h$ with $h \ge 1$, the bare ground state energy evaluates to:
\begin{equation}
E_\Omega = -\frac{h(h-1)}{2}.
\end{equation}
\end{lemma}

\begin{proof}
By definition, $S_\Omega = \{-h+1, \dots, 0\}$. Summing the integers:
\[
E_\Omega = \sum_{j=-h+1}^0 j = - \sum_{m=1}^{h-1} m = - \frac{(h-1) h}{2} = -\frac{h(h-1)}{2}.
\]
\end{proof}

\begin{theorem}[Spectra and Commutation Relations]
\label{thm:h0_commutators}
\begin{enumerate}
    \item On any occupation basis state $|S\rangle$ ($S \subseteq \Lambda^*$):
    \begin{equation}
    H_0 |S\rangle = \left( \sum_{k \in S} k \right) |S\rangle, \qquad \hat{P} |S\rangle = \left( \sum_{k \in S} k - E_\Omega \right) |S\rangle.
    \end{equation}
    \item For any mode $p \in \Lambda^*$, the creation and annihilation operators satisfy:
    \begin{align}
    [H_0, c_p^\dagger] &= [\hat{P}, c_p^\dagger] = p\, c_p^\dagger, \\
    [H_0, c_p] &= [\hat{P}, c_p] = -p\, c_p.
    \end{align}
\end{enumerate}
\end{theorem}

\begin{proof}
We prove each statement:
\begin{enumerate}
    \item Each $n_k$ is diagonal on basis states: $n_k |S\rangle = |S\rangle$ if $k \in S$, and $0$ otherwise. Summing over $k \in \Lambda^*$ gives $H_0 |S\rangle = (\sum_{k \in S} k) |S\rangle$. Subtracting $E_\Omega I$ gives the eigenvalue for $\hat{P}$.
    \item Using the bilinear CAR commutator $[n_k, c_p^\dagger] = \delta_{k, p} c_p^\dagger$ from Theorem~\ref{thm:algebraic_car}(2):
    \[
    [H_0, c_p^\dagger] = \left[ \sum_{k \in \Lambda^*} k\, n_k, \; c_p^\dagger \right] = \sum_{k \in \Lambda^*} k\, [n_k, c_p^\dagger] = \sum_{k \in \Lambda^*} k (\delta_{k, p} c_p^\dagger) = p\, c_p^\dagger.
    \]
    Similarly, $[n_k, c_p] = -\delta_{k, p} c_p$ yields $[H_0, c_p] = -p\, c_p$.
    Because $\hat{P} = H_0 - E_\Omega I$ differs from $H_0$ only by a central scalar multiple of the identity, the commutators with $\hat{P}$ are identical.
\end{enumerate}
\end{proof}

Theorem~\ref{thm:h0_commutators} confirms that $c_p^\dagger$ raises the total physical momentum energy by $p$, while $c_p$ lowers it by $p$. This provides the kinematic foundation for constructing particle-hole excitations and density modes in the subsequent chapters.
